% Werke X/1 (1917), native177–178, printed170–171.
\begingroup
\setlength{\headheight}{25pt}
\SourceLeaf{177}\ArithmeticHead{170}{L}{ANALYSIS.\quad NACHLASS.}
\NumberedItem{6}
\HistoricCitation{[Zettel in Fh Nr. 2, Kapsel 50.]}
\par\indent Ponendo \(\sin\text{ lemn }\varphi=\sin\text{ circ. }\chi\) erit
\[\begin{array}{r@{\;}r@{\;}l}
\chi=\frac{\pi}{\varpi}\varphi-{}&0{,}086266&\sin2\varphi\\
+{}&1866&\sin4\varphi\\
-{}&53&\sin6\varphi\\
+{}&2&\sin8\varphi
\end{array}\]
\noindent [Auf der Rückseite]
\[Px=\frac{x-\frac1{60}x^5-\frac1{10080}x^9\ldots}{e^{\frac12\frac{xx}{\varpi\varpi}\pi}}\]
\[Q(1+\surd-1)\pi=4e^{\frac12\varpi}.\]
\begin{center}\small BEMERKUNGEN.\end{center}
\renewcommand{\thefootnote}{*)}
\par\indent Die in den artt. [1.] bis [5.] wiedergegebenen Leisteaufzeichnungen gehören zu den \textls[120]{irregulären}\footnote{Während die in den Abschnitten [I.], [II.], [III.] wiedergegebenen Leisteaufzeichnungen in fortlaufender Reihenfolge gemacht sind und den Tagebuchnotizen des Jahres 1797 entsprechen, sind die in diesem Abschnitt [IV.] zusammengestellten Notizen auf leer gebliebenen Stellen von Durchschußblättern, die \textsc{Gauss} früher mit Aufzeichnungen andern Inhalts beschrieben hatte, eingetragen. Wir bezeichnen solche Leistenotizen im Text und weiterhin als \textls[120]{irreguläre}. Vergl. den Artikel 2 des Aufsatzes »Über \textsc{Gauss}' Arbeiten zur Funktionentheorie«.}.
In [1.], [2.] sind \(P,Q\) die früher mit \(M,N\) bezeichneten Zähler und Nenner des sinus lemniscaticus; diese Bezeichnungen hat \textsc{Gauss} auch auf dem Werke III, S. 404---406, art. [1.]---[4.] abgedruckten Zettel und in sehr viel späterer Zeit in den Handbüchern 16, Bb und 19, Be (siehe Werke III, S. 408---412) benutzt. Zu den Formeln in [1.] vergleiche man Werke III, S. 411, 412; in den Koeffizienten von \(\frac{P^8}{Q^8}\) sind in den Zählern die Potenzen von \(n\) zu ergänzen. Die Aufzeichnung [3.] erscheint schon äußerlich durch das Zeichen \(\varpi\) für die halbe lemniskatische Periode als irregulär; der hingeschriebene Quotient ist die von Juli 1798 ab in der Scheda Aa\renewcommand{\thefootnote}{**)}\footnote{Die Scheda Aa ist begonnen Juli 1798; ihr Inhalt hängt mit den Tagebuchnotizen Nr. 92, 94, 95 (Juli---Oktober 1798) zusammen. Sie dürfte Oktober 1798 abgeschlossen worden sein, denn im November 1798 beginnt die Scheda Ab.} (siehe Werke III, S. 416 ff.) ständig mit \(P\) bezeichnete Funktion, die die Eigenschaft
\[P(\varpi+x)=-P(x)\]
\clearpage
\SourceLeaf{178}\fancyhead{}\fancyhead[R]{\small171}\fancyhead[C]{\scriptsize\shortstack{ÄLTESTE UNTERSUCHUNGEN ÜBER\\ LEMNISKATISCHE FUNKTIONEN.\quad IV.}}
\noindent besitzt und sich als die \textsc{Jacobi}sche Thetafunktion \(H(x)\)\renewcommand{\thefootnote}{*)}\footnote{Vergl. etwa \textsc{Jacobis} Werke I, S. 224.} für den lemniskatischen Fall charakterisieren läßt, während die Funktionen \(M,N\) den \textsc{Weierstrass}schen Al-Funktionen\renewcommand{\thefootnote}{**)}\footnote{Vergl. \textsc{Weierstrass}' Werke I, S. 6.\ClosingRule} entsprechen. Derselbe Quotient findet sich auch auf dem in [6.] abgedruckten Zettel. Die daselbst und in [4.] auftretenden Reihenentwickelungen nach den Sinus der Vielfachen eines Winkels weisen auf die Tagebuchnotiz Nr. 91 b vom Juli 1798 hin. Die in [4.] in dem ersten Ansatz für \(\sin\text{ lemn }a\) als Koeffizient von \(\sin a\) auftretende Zahl stand ursprünglich auch ebenso (mit 60 in der 6. und 7. Dezimale) im \textit{Tagebuch}; die dann folgende Entwickelung in [4.] ist viel vollständiger und in den 7. und 8. Dezimalen genauer als die des \textit{Tagebuchs}, sie ist also jedenfalls später, \textls[120]{und damit ist dieser ganze Komplex von Aufzeichnungen auf Juli 1798, kurz vor Beginn der Scheda, Aa datiert.}

Das am Schluß von [4.] hingeschriebene Integral stellt den Zusammenhang mit der Aufzeichnung [5.] her, die eine geschicktere und schärfere Ausführung der in [I.] [8.] (oben S. 149 dieses Bandes) begonnenen Rechnung enthält. In [5.] bedeutet \(\Pi\) wieder die Viertelperiode
\[\Pi=\int_0^1\frac{dp}{\sqrt{1-p^4}}=\frac{\varpi}{2};\]
setzt man in diesem Integral \(p=\sqrt{1-\xi^2}\), so erhält man die Form
\[\Pi=\frac1{\sqrt2}\int_0^1\frac{d\xi}{\sqrt{(1-\xi^2)(1-\frac12\xi^2)}},\]
aus der die erste in [5.] angegebene Reihe für \(\Pi\) hervorgeht. Die Reihe für
\[\sqrt{\frac32}.\frac{\varpi}{\pi}=\sqrt6\frac{1{,}3110\ldots}{3{,}1415\ldots}\]
findet sich auch in der Tagebuchnotiz Nr. 91 a (Juli 1798). Die mit Hilfe dieser Reihe ausgeführte numerische Rechnung liefert einen Wert für \(\frac{\varpi}{\pi}\), der zwar viel genauer ist als der in [I.] [8.] angegebene, jedoch wird auf S. 24 der Scheda Aa ein noch genauerer Wert berechnet (es folgen nämlich auf die vierzehnte Dezimalstelle 3 noch die Stellen 164), so daß also die hier gegebene Rechnung die frühere ist. Damit ist auch [5.] auf Juli 1798 datiert, wobei allerdings noch die Möglichkeit offen bleibt, daß die erste Reihe für \(\Pi\) schon vor den übrigen Notizen hingeschrieben worden ist, zu einer Zeit, als \textsc{Gauss} noch \(\Pi\) und nicht \(\frac{\varpi}{2}\) als Zeichen für die Viertelperiode benutzte.
\begin{flushright}\textsc{Schlesinger.}\end{flushright}
\clearpage
\endgroup
